\begin{abstract}
The values of causal effects and counterfactuals are typically only partially identified from observational data in the presence of unobserved confounding. With discretely-valued variables, under certain conditions, sharp bounds can be derived by solving one linear program. This is then computationally tractable for many causal questions. In this paper, for every diagram in which each confounded component has at most two hidden causes, we give necessary and sufficient conditions for the sharp bounds to be equal to the optimum of one linear program for every (positive) input probability table. This algorithm is implemented in a new Python package, called \pid{}, that allows users to compose assumptions, queries, and data to obtain sharp bounds in a computationally tractable manner.
\end{abstract}

\noindent\textit{Keywords:}
partial identification;
causal inference;
linear programming;
sensitivity analysis;
instrumental variables;
Python.
